\documentclass{article}
\usepackage{amsmath,amsthm,enumitem,listings}
\newtheorem{definition}{Definition}
\newtheorem{example}{Example}
\newcommand{\Oh}[1]{\mathcal{O}\!\left(#1\right)}
\newenvironment{keypoint}{\par\noindent\textbf{Key point.}\ }{\par}
\title{Notes on Sorting}
\author{Ruth Ng}
\begin{document}
\maketitle
\tableofcontents
\section{Comparison sorts}
\begin{definition}[Comparison sort]
A sort that orders its input only by comparing pairs of elements.
\end{definition}
\begin{keypoint}Any comparison sort needs $\Omega(n\log n)$ comparisons in the worst case.\end{keypoint}
\subsection{Merge sort}
Merge sort runs in $\Oh{n\log n}$ time and uses $\Oh{n}$ extra space.
\begin{lstlisting}[language=Python]
def merge_sort(xs):
    if len(xs) <= 1:
        return xs
    mid = len(xs) // 2
    return merge(merge_sort(xs[:mid]), merge_sort(xs[mid:]))
\end{lstlisting}
\begin{example}
Sorting $[3,1,2]$ takes three comparisons.
\end{example}
\section{Counting sorts}
\begin{description}
  \item[Counting sort] stable, $\Oh{n+k}$.
  \item[Radix sort] stable, $\Oh{d(n+k)}$.
\end{description}
\begin{tabular}{|l|r|r|}
\hline
Algorithm & Time & Space \\ \hline
Merge & $n\log n$ & $n$ \\ \hline
Counting & $n+k$ & $k$ \\ \hline
\end{tabular}
\end{document}
